\chapter{Discussion}\label{ch:discussion}

\section{What the work shows}

One guess about $\pi$ led, through one algorithm, to a surprising amount of mathematics and music. The guess $3 + k/(7k+1)$ is a continued fraction; a continued fraction is a rhythm whose levels are its terms (\cref{thm:nested}); the same algorithm builds tunings, Euclidean rhythms, mode-locking staircases and a factoring method; coprimality, which every level of the rhythm obeys, is where $6/\pi^2$ enters; and the self-similarity of the continued fraction is the common root of a family of classical fractals. Along the way we found the precise meaning of two numbers that the site's sources disagreed about (32,996 and 33,102), and we checked every numerical claim by script.

The two research threads gave concrete answers. For prime convergents, a heuristic built from Lévy's theorem and the equidistribution of convergents modulo $m$ predicts $2\ln 2 \cdot \ln N$ prime denominators among the first $N$ convergents of a typical number, with the correction factor turning out to be $\zeta(2)$; a sample of random reals fits the predicted growth, and $\pi$ is typical. For the Montgomery--Dyson connection, one set of statistics places a crystal, $\pi$'s rhythm, the zeta zeros, random matrices, the primes and Poisson points on a single scale, and the explicit formula makes Dyson's quasicrystal remark visible. Neither thread proves a theorem about $\pi$ or the zeta function, and we have been careful to say so.

\section{Limitations}

\begin{itemize}
\item \emph{Heuristics are not theorems.} \Cref{heur:primeconv} treats primality of different denominators as independent events; the agreement with data supports it but does not prove it. Chapter~\ref{ch:openq} proves it in expectation, not for almost every number. \addedon{\revA} Chapter~\ref{ch:openq2} proves it for almost every number; for $\pi$ itself it is still a heuristic. \addedon{\revB}
\item \emph{Finite samples.} The prime-convergent test uses \RandM{} random reals and $1{,}000$ convergents each; the slope fit has a standard error of $\RandSlopeSE$ and is slightly above the prediction because of early terms. A larger sample, and longer expansions, would sharpen the test.
\item \emph{Low zeros.} We used the first \NZeros{} zeros, which we computed ourselves. At these heights finite-size corrections are visible \citep{Bogomolny2006,Berry1988}; Odlyzko's tables of zeros at much greater heights \citep{Odlyzko1987} would show closer agreement with the GUE.
\item \emph{Probable primes.} Large ``primes'' are Baillie--PSW probable primes. No composite is known to pass that test, but it is not a proof of primality.
\item \emph{Score facts.} The facts about the Suite come from the site's data, which its builders read from the engraved score; we checked the bar arithmetic but not each note.
\item \emph{Literature search.} The audit of Chapter~\ref{ch:new} is as careful as we could make it, but the music-theory literature in particular is large, and related work may exist that we did not find.
\end{itemize}

\section{Open questions}\label{sec:open}

Six questions were left open when this work was first written. Chapter~\ref{ch:openq} reports what further work on them produced; the status after that work is given in brackets, and later updates follow it. \addedon{\revA}
\begin{enumerate}
\item Does $\pi$ have infinitely many prime convergent denominators? Infinitely many prime numerators? [Open. We tested $q_n$ and $p_n$ to $n = \OqPiN$ using \OqNTm{} million certified terms; the counts stay inside the heuristic's 95\% range (Section~\ref{sec:oq1}).] \addedon{\revA} [Update: the period of $e$'s denominators modulo a prime $\ell$ divides $6\ell$, now proved (\cref{prop:ob-e}).] \addedon{\revB} [Update: numerators and denominators searched to $n = 50{,}000$. Four new probable-prime denominators, at $n = 25{,}332$, $27{,}651$, $32{,}394$ and $37{,}721$, and a new probable-prime numerator, at $n = 42{,}037$, all found by Teddy Rook. They form a cluster, and the prime convergents lean toward one note; both look like chance, and neither says where the next one is (Sections~\ref{sec:th1} and~\ref{sec:th6}). For $e$, the period modulo $\ell$ is exactly $6\ell$ and prime denominators can occur only at $n = 2$ or $n \equiv 1 \pmod 3$ (\cref{prop:th-flip,thm:th-div}).] \addedon{\revC}
\item Are there infinitely many Fibonacci primes? Pell primes? [Open. The known lists agree reasonably with the Grantham--Granville heuristic (Section~\ref{sec:oq2}).] \addedon{\revA}
\item Can \cref{heur:primeconv} be proved in a model? For example, for almost every $x$, is the number of $n \le N$ with $q_n$ prime asymptotic to $2\ln 2 \cdot \ln N$? Even an upper bound of that order for almost every $x$ would be interesting. [Partly answered. For a random $x$ the expected count is $2\ln 2\cdot\ln N + O(1)$, a theorem (\cref{cor:oq-N}); for almost every $x$ the count lies between $\ln N$ and $2\ln N$ asymptotically, using recent work on the Duffin--Schaeffer conjecture (\cref{thm:oq-ae}). The almost-everywhere constant $2\ln 2$ is still a conjecture.] \addedon{\revA} [Update: answered for almost every $x$. The count is $2\ln 2\cdot\ln N\,(1 + o(1))$ for almost every $x$, with an explicit error term (\cref{thm:ob-ae}). Whether $\pi$ is one of these numbers is open.] \addedon{\revB} [Update: for $\pi$ itself, its prime convergents to $n = 50{,}000$ now favor the constant $2\ln 2$ over the value without the factor $\zeta(2)$ by about \RaFiLR{} to 1 (evidence, Section~\ref{sec:th7}).] \addedon{\revC}
\item Are $\pi$'s terms unbounded? Do they obey the Gauss--Kuzmin law? [Open. Calibrated tests on \OqNTm{} million terms are all consistent with typical behavior (Section~\ref{sec:oq4}).] \addedon{\revA} [Update: two more tests, the Doeblin--Lenstra law and the serial correlation of the terms, are also consistent (Section~\ref{sec:ob4}).] \addedon{\revB} [Update: every prime below $1{,}000$ divides $\pi$'s first two million denominators at its typical rate $1/(\ell + 1)$, and the share of denominators with no prime factor below $10^7$ matches the typical-number prediction (Section~\ref{sec:th7}).] \addedon{\revC}
\item How does the spectrum of the arithmetic correlation matrix of Section~\ref{sec:corr} compare with random-matrix benchmarks, and does $6/\pi^2$ appear in its bulk? [Partly answered. Its eigenvalues have mean 1 and variance tending to 2, and its determinant and geometric mean are explicit (proved); $6/\pi^2$ appears in a lower bound $(2/\pi^2)\ln^2 m$ for the top eigenvalue, not in the first two moments; its spacings look like the GOE's (Section~\ref{sec:oq5}).] \addedon{\revA} [Update: the top eigenvalue grows faster than any power of $\ln m$ and at most like $m^{o(1)}$ (\cref{thm:ob-top}); the spectrum contains near-copies of the small gcd matrices, one family for each value of $\lfloor m/p\rfloor$ (\cref{prop:ob-copies}), and with these set aside the spacing ratio falls from the GOE value toward the middle, so the GOE-like look was partly an artifact.] \addedon{\revB} [Update: the extremes are settled to leading order: $\ln\lambda_{\max}$ and $\ln(1/\lambda_{\min})$ are both $(2 + o(1))\sqrt{\ln m/\ln\ln m}$ (\cref{thm:th-top,thm:th-min}). The eigenvalue distribution converges to a limiting law (\cref{thm:th-law}), whose known exact constants do not involve $6/\pi^2$. The shape of that law is open.] \addedon{\revC} [Update: the inverse has the sign pattern of Liouville's function, so both extreme eigenvalues are Perron eigenvalues of non-negative matrices built from common multiples (\cref{prop:18mirror,prop:18harmonic}). The ratio of $1/\lambda_{\min}$ to $\lambda_{\max}$ is about $2.6$ at $m = 3\times10^7$ and near $2.57$ for the gcd matrix; whether it has a limit is open (Section~\ref{sec:18mirror}).] \addedon{\revD}
\item Rank all equal temperaments up to 200 notes by how well they approximate several just intervals at once. Do prime sizes stand out? [Answered: no, by eight measures and by an equidistribution theorem (Section~\ref{sec:oq6}).] \addedon{\revA} [Update: why 171 wins: it is the equal temperament fixed by three tiny commas (Section~\ref{sec:ob6}), a fact known to tuning theorists.] \addedon{\revB}
\end{enumerate}

\section{Future work}

On the mathematical side: find the second-moment estimate that would prove \cref{heur:primeconv} for almost every $x$ (Section~\ref{sec:oq3}); find the exact growth of the top eigenvalue of the correlation matrix; and repeat the spacing statistics with zeros at large height. \addedon{\revA} [Update: the second-moment estimate is now proved (\cref{thm:ob-ae}); the top eigenvalue's exact growth is still open, between the two bounds of \cref{thm:ob-top}.] \addedon{\revB} [Update: the growth of the top eigenvalue is now known to leading order (\cref{thm:th-top}); a second-order term, and a description of the limiting bulk law, are the open problems.] \addedon{\revC} [Update: on the computational side, certificates of primality for the five new probable primes of Section~\ref{sec:th1}, and a search of $\pi$'s convergents beyond $n = 50{,}000$, are the next steps.] \addedon{\revC} On the musical side: the method of Chapter~\ref{ch:suite} invites pieces in tunings whose number of notes is itself a convergent denominator (a pi movement in a 113-note tuning, for instance), and pieces built from numbers with chosen terms, whose form a composer could design by picking the continued fraction. The companion site's instruments already do much of this; the score of each new piece is a continued fraction made audible.

